%% Dans ce fichier, les �ventuels accents sont cod�s
%% en ISO-8859-15 (aka latin9)
%% La ligne suivante donne trois exemples 
%% e accent aigu : �, e accent grave : �, e-dans-l'o : �
%%%%%ENTREES%%%%
\begin{tabarg}
\multicolumn{3}{l}{\titretab {Entr\'ees}}\\
\hline
$X$ & Vecteur, \texttt{float} & Param�tre avec valeurs � remplacer (unit� quelconque)\\
$Y$ & Vecteur, \texttt{float} & Param�tre selon lequel s'effectue la s�lection (unit� quelconque)\\
$Hors\_int$ & Coeff, \texttt{float} & 0 ou 1 pour sp�cifier si les valeurs � remplacer sont dans ou hors de l'intervalle [$Crit\_{min}$, $Crit\_{max}$] \\
$Crit\_{min}$ & Coeff., \texttt{float} & valeur mini de l'intervalle de test de  $Y$\\
$Crit\_{max}$ & Coeff., \texttt{float} & valeur maxi de l'intervalle de test de $Y$\\
$Y_{mnq}$ & Coeff., \texttt{float} & valeur manquante de $Y$ \\
$ValR$ & Coeff., \texttt{float} & valeur de remplacement de $X$
soit \texttt{\$Arg1.misval}\\

\end{tabarg}

\medskip
%%%%%SORTIES%%%%
\begin{tabarg}
\multicolumn{3}{l}{\titretab {Sorties}}\\
\hline
$X_{rpl}$ & Vecteur, \texttt{float} & $X$ avec valeurs remplac�es (\$Arg1.unit)\\
\end{tabarg}
